% FILE: 07_open_questions_and_next_work.tex
% Project: experiments-visualizer
% Thesis Role: process-intelligence-simulation-dashboard

\section{Open Questions and Next Work}
\label{sec:experiments-visualizer-openq}

\subsection{Known Gaps}
\label{sec:experiments-visualizer-openq-gaps}

Four confirmed gaps exist in the \texttt{experiments-visualizer} as of the current evidence
state. These gaps are documented from direct inspection of the source artifacts and are not
inferences.

\textbf{Gap 1: No automated runtime test harness.}
\texttt{visualizer-validation.ts} exercises TypeScript type shapes at compile time only. There
is no runtime test harness, Jest suite, or browser automation framework confirming that the
MAPE-K loop executes correctly end-to-end, that the SHA-256 chain grows correctly across
multiple blocks, that the EWMA drift detector correctly triggers UCL/LCL violations, or that
the DECLARE validator correctly evaluates streaming traces. The v30.1.1 conformance claim in
\texttt{LIVESTREAM\_STANDARDS.md} is a specification document, not a test result.

\textbf{Gap 2: No visualizer-specific ALIVE checkpoint.}
The \texttt{PROCESS\_INTELLIGENCE\_ALIVE\_001.md} verdict covers the overall process
intelligence program. No checkpoint document exists that issues an ALIVE verdict specifically
for the \texttt{experiments-visualizer} dashboard. The project's ALIVE status is inferred
from the broader checkpoint, making the attribution PARTIAL under the thesis immutability
doctrine.

\textbf{Gap 3: Four independent SHA-256 implementations.}
SHA-256 is independently implemented in \texttt{blockchain.js}
(\texttt{CryptographicAuditChain.sha256}), \texttt{autonomic.js} (\texttt{sha256}),
\texttt{ledger.js} (\texttt{sha256}), and \texttt{dashboard.js} (\texttt{sha256}). It is
unconfirmed whether these are intentionally independent (defense-in-depth) or should reference
a single canonical implementation. If the implementations differ, blocks manufactured by
different modules may produce inconsistent hashes for identical inputs, undermining the
tamper-evidence guarantee.

\textbf{Gap 4: AGI swarm specification without implementation.}
\texttt{LIVESTREAM\_STANDARDS.md} specifies a 5-agent AGI swarm architecture (Stream Director,
Telemetry Auditor, A* Alignment Referee, Drift Sentry, Ledger Custodian). No agent
orchestration code, OBS bindings, or inter-agent communication protocol is present in the
codebase. The swarm architecture is a specification document only; its implementation status
is unconfirmed.

\subsection{Deferred Gates}
\label{sec:experiments-visualizer-openq-deferred}

The following proof gates are deferred pending further work:

\begin{enumerate}
  \item \textbf{Runtime execution gate:} Automated confirmation that all seven subsystems
        execute correctly in a live browser session, producing mutually consistent outputs
        without error. Requires a browser automation framework (e.g.\ Playwright or Puppeteer)
        wired to the dashboard entry point.
  \item \textbf{Visualizer-specific ALIVE checkpoint:} A checkpoint document issued specifically
        for the \texttt{experiments-visualizer} satisfying all gate criteria independently of
        the broader process intelligence program checkpoint.
  \item \textbf{WASM provenance tracing:} Confirmation that \texttt{bindings.d.ts} is the
        authoritative projection of the Rust wasm4pm WASM module's type surface, with a
        traceable connection to a Cargo build script, \texttt{wasm-bindgen} configuration, or
        equivalent generation mechanism. Currently the file is labeled ``generated compile-time
        TypeScript type definitions'' but no generation pipeline is visible in the visualizer
        directory.
  \item \textbf{SHA-256 canonicalization:} Confirmation that the four SHA-256 implementations
        produce identical output for identical inputs, or formal documentation of why
        independent implementations are required.
\end{enumerate}

\subsection{Research Risks}
\label{sec:experiments-visualizer-openq-risks}

\textbf{Risk 1: Specification-implementation gap.}
The \texttt{LIVESTREAM\_STANDARDS.md} v30.1.1 conformance claim is the most aggressive quality
assertion in the codebase, prohibiting all mocks, placeholders, and stubs. Without automated
runtime testing, this claim cannot be verified. If any of the 13 JS modules contains a
stubbed path that is activated only under specific runtime conditions, the conformance claim
is falsified without any existing test being able to detect it.

\textbf{Risk 2: Single-threaded execution limits.}
The browser's single-threaded JavaScript event loop imposes a fundamental constraint on the
MAPE-K tick cycle: all seven subsystems must complete their computation within a single event
loop iteration without blocking the rendering thread. If the A* alignment solver's
priority-queue search takes too long on a complex trace, the dashboard will appear to freeze.
No performance benchmarks or tick-budget measurements are present in the codebase.

\textbf{Risk 3: Indirect WASM boundary.}
The connection between \texttt{bindings.d.ts} and the actual Rust wasm4pm WASM module is not
directly traceable from the visualizer directory. If the TypeScript declarations have drifted
from the actual WASM exports, the compile-time type gate passes while the runtime boundary
produces type errors. This risk cannot be resolved without tracing the provenance of
\texttt{bindings.d.ts} to its upstream source.

\subsection{Next Receipted Work}
\label{sec:experiments-visualizer-openq-next}

The following work items, if completed, would close the deferred gates and advance the project
from PARTIAL to ALIVE status. Each item must be receipted upon completion under the process
intelligence receipt chain protocol.

\begin{enumerate}
  \item Manufacture a browser automation test suite (Playwright recommended) covering the
        full MAPE-K tick cycle, SHA-256 chain growth, DECLARE constraint evaluation over
        streaming traces, and EWMA UCL/LCL breach detection. Receipt the test results as
        a \texttt{ReceiptShapeTs}-conforming record.

  \item Issue a visualizer-specific ALIVE checkpoint in
        \texttt{/Users/sac/process-intelligence/checkpoints/} satisfying all four gates
        independently of the broader program checkpoint.

  \item Canonicalize the SHA-256 implementation into a single shared module imported by
        \texttt{blockchain.js}, \texttt{autonomic.js}, \texttt{ledger.js}, and
        \texttt{dashboard.js}. Alternatively, document the rationale for independent
        implementations and confirm output equivalence.

  \item Trace the provenance of \texttt{bindings.d.ts} to its upstream source in the Rust
        wasm4pm build pipeline and document the traceability chain in a receipt.

  \item Implement or formally specify the 5-agent swarm roles defined in
        \texttt{LIVESTREAM\_STANDARDS.md} as either automated test agents or documented
        human observer protocols with formal handoff criteria.
\end{enumerate}
